% !TEX root = ../../main.tex

\section{本章小结}

本章围绕 HearBridge 系统中的手势识别服务实现、句子视频识别流程与实验结果分析展开说明。

首先，
针对移动端实时训练与识别场景，
本章介绍了自采样本采集链路、手势特征构造、词级模型训练、模型版本发布和实时识别闭环。

在该流程中，
移动端负责采集用户手势图像帧并通过 WebSocket 发送至 Python 识别服务，
Python 端完成 MediaPipe 关键点检测、特征提取、样本保存、模型训练和实时预测，
后台管理端则负责样本管理、模型训练触发、模型版本发布和发布模型重载。

该部分验证了系统能够完成从移动端样本采集、服务端特征转换、模型训练发布到移动端实时识别展示的完整闭环。

其次，
针对句子视频识别场景，
本章说明了基于 WLASL 小词表数据的词级分类模型构建过程。

系统采用 $20\times232$ 的 plus 特征作为词级分类输入，
并在 25 类小词表上完成模型训练与测试。

实验结果表明，
词级模型虽然在 Top-1 分类准确率上仍有提升空间，
但 Top-3 准确率较高，
能够为后续词段 Top-K 候选保留、候选序列生成和语义修正提供有效基础。

在句子视频识别流程方面，
本章进一步介绍了上传视频后的帧级检测缓存、滑动窗口输入生成、逐窗口词级分类、词段合并、NMS 抑制、Gloss 序列生成和 DeepSeek 语义修正机制。

该流程将长度可变的短句视频转换为多个固定长度候选窗口，
再通过词级模型得到窗口预测结果，
最终经词段后处理生成原始 Gloss 序列。

随后，
后端基于词段 Top-K 候选生成有限候选序列，
并约束 DeepSeek 只能从候选集合中选择一个候选编号，
从而在保持视觉证据约束的前提下完成语义修正和中文自然化表达。

在实验分析部分，
本章对 core61 拼接句子视频测试集进行了多随机种子实验。

结果表明，
原始句子级识别已经能够完成部分短句的完整匹配；
候选集合理论上限明显高于原始识别结果，
说明 Top-K 候选保留和候选序列生成具有较大的可恢复空间；
DeepSeek 候选重排序相比原始结果具有稳定但有限的平均提升，
说明受约束语义修正在当前小词表条件下具有一定潜力。

同时，
实验也表明当前系统仍存在插入误报、漏词、冗余输出、局部替换错误和候选排序未命中等问题。

这些问题一方面来自滑动窗口机制对动作边界的间接建模，
另一方面也与单词级训练资源不足、词表规模较小和相似动作类别混淆有关。

因此，
当前方案更适合作为有限词表短句识别和系统功能验证方案，
尚不能等同于开放域连续手语识别模型。

最后，
本章还对更直接的连续手语识别方案进行了探索与说明。

相比词级分类与滑动窗口方案，
基于 Transformer、Conformer 和 CTC 的连续识别方法在任务形式上更接近真实连续手语识别，
能够直接从长视频特征序列中预测 Gloss 序列。

但是，
该方向对连续数据质量、序列建模能力、训练策略、参数调优和解码方法要求更高。

在当前毕业设计周期和系统交付目标下，
相关实验尚未达到可稳定接入系统的程度。

因此，
本项目最终采用训练难度更可控、工程链路更稳定的词级分类与滑动窗口方案，
优先完成可运行、可测试、可展示的端到端识别系统。

综上，
本章完成了 HearBridge 手势识别服务从移动端实时识别到句子视频识别的核心实现与实验验证。

当前系统已经能够支撑样本采集、模型训练、模型发布、实时识别、句子视频上传识别和中文语义表达等主要功能。

后续工作可从扩充高质量手语视频资源、扩大词表规模、优化词级模型特征表达、改进词段边界处理、增强候选排序能力以及继续探索连续识别模型等方面进一步提升系统性能。